\documentclass{article}
\title{Hedgerow Width and Ground Beetle Species Richness on Arable Farms}
\date{}
\begin{document}
\maketitle

\begin{abstract}
We tested whether wider hedgerows hold more ground beetle species on arable farms. Across 24 plots on 12 farms, species richness (species per plot) rose by 3.1 species per additional metre of hedgerow width. A second survey measured beetle activity as captures per month and found peak activity in June.
\end{abstract}

\noindent Keywords: ground beetles, hedgerows, species richness, agroecology

\section{Introduction}

Hedgerows shelter ground beetles that eat crop pests and weed seeds (Lang, 2016). Farms in eastern England removed 22 percent of their hedgerows between 1990 and 2010 (Weber, 2015).

The present survey related species richness to hedgerow width. The hypothesis was that wider hedgerows hold more species.

\section{Richness Method}

A total of 24 plots were sampled on 12 farms in an agricultural landscape of 40 square kilometres. Plots lay along hedgerows 1 to 7 m wide. Each plot held one transect.

Pitfall traps were placed 1 m apart along a 20 m transect. Each trap held propylene glycol, which killed the beetles within 1 hour and preserved them for counting. The traps were emptied every 7 days from April to September.

Soil moisture was sampled 6 inches below the surface.

The researchers identified each beetle to species.

We kept the plots at least 100 m apart so that the samples stayed independent. Because all plots lay in one region, the samples were independent within that region only, a limitation of the richness survey.

Each plot's hedgerow was described by its width and its age. Width was the mean of measurements at 5 points per plot. Age came from the farm records.

Beetles were counted in the laboratory under a microscope by two observers who worked independently. Their counts agreed for 97 percent of samples, and the counts were used without correction.

\section{Richness Results}

Species richness rose by 3.1 species per additional metre of hedgerow width (95% CI [2.2, 4.0]). Mean species richness was 12.3 species per plot (SD = 2.1). Because the interval excluded zero, width and species richness were positively associated.

Plots wider than 5 m held 16 species on average, 3.7 species above the mean of all plots. That excess agreed with the positive association between width and species richness.

Three species occurred only in hedgerows wider than 6 m, so width may have decided which species could persist.

Species richness did not depend on farm size (r = 0.12, p = .58). Farm size was an unlikely source of the width effect.

Hedgerows wider than 5 m held 9 forest specialists, species that live mainly in woodland, and hedgerows narrower than 3 m held 2. The difference may have reflected the shade that the wider hedgerows offered these species.

Hedgerow age rose with width (r = 0.64, p = .001), so the older hedgerows were also the wider ones.

\section{Activity Method}

A second survey recorded when ground beetles were most active.

Beetles were sampled along hedgerows.

Species were identified from published keys.

Only one season was sampled, and the pattern may differ in other years.

\section{Activity Results}

Beetle activity peaked in June, with 41 percent of all captures. Activity fell below 5 percent per month after August. Beetles were therefore most active in early summer.

We were delighted to find Carabus auratus in the third hedgerow.

The timing of activity showed no link to hedgerow width (F = 0.8, p = .45). Width thus appeared to shape how many species a plot held and to leave their season of activity unchanged.

\section{Discussion}

Wider hedgerows hold more species, which suggests that width adds habitat. The result extends earlier work on hedgerows as shelter for beetles (Lang, 2016) to the width of the hedgerow.

The pattern matches earlier surveys in Denmark (Holm, 2012).

Older hedgerows are wider in this sample, so age may explain part of the gain. A survey of hedgerow age on the same plots would separate the two effects.

A 5 m strip removes land from the crop; the same strip holds 16 species on average. Farmers therefore weigh a loss of cropland against a gain in species.

The gain relates to width, not to hedgerow length, since every transect was 20 m long.

We suggest that agri-environment schemes consider payments for hedgerows wider than 5 m.

\section{Conclusion}

Hedgerows wider than 5 m may hold more ground beetle species on arable farms. Future work should follow the same plots across at least three seasons.


\end{document}
